% Part: first-order-logic
% Chapter: axiomatic-deduction
% Section: proving-things-quant

\documentclass[../../../include/open-logic-section]{subfiles}

\begin{document}

\olfileid[id]{fol}{axd}{prq}

\olsection{\usetoken{P}{derivation} dengan Kuantor}

\begin{ex}
Mari kita berikan !!a{derivation} atas $(\lforall[x][!A(x)] \land
\lforall[y][!B(y)]) \lif \lforall[x][(!A(x) \land !B(x))]$.

Pertama, perhatikan bahwa
\begin{align*}
  (\lforall[x][!A(x)] \land \lforall[y][!B(y)]) & \lif \lforall[x][!A(x)]\\
  \intertext{merupakan instansiasi dari \olref[prp]{ax:land1}, sedangkan}
  \lforall[x][!A(x)] & \lif !A(a) \\
  \intertext{merupakan instansiasi dari \olref[qua]{ax:q1}. Jadi, menurut \olref[pro]{prop:chain}, kita tahu bahwa}
  (\lforall[x][!A(x)] \land \lforall[y][!B(y)]) & \lif !A(a) \\
  \intertext{!!{derivable}. Demikian pula, karena}
  (\lforall[x][!A(x)] \land \lforall[y][!B(y)]) & \lif \lforall[y][!B(y)] \qquad\text{dan}\\
    \lforall[y][!B(y)] & \lif !B(a)\\
    \intertext{masing-masing merupakan instansiasi dari \olref[prp]{ax:land2} dan \olref[qua]{ax:q1},}
    (\lforall[x][!A(x)] \land \lforall[y][!B(y)]) & \lif !B(a)\\
    \intertext{dapat diturunkan menurut \olref[pro]{prop:chain}. Dengan menggunakan instansiasi yang sesuai dari \olref[prp]{ax:land3} dan dua penerapan~\MP, kita melihat bahwa}
    (\lforall[x][!A(x)] \land \lforall[y][!B(y)]) & \lif (!A(a) \land !B(a))\\
    \intertext{dapat diturunkan. Sekarang kita dapat menerapkan \QR{} untuk memperoleh}
(\lforall[x][!A(x)] \land \lforall[y][!B(y)]) & \lif \lforall[x][(!A(x) \land !B(x))].
\end{align*}
\end{ex}

\end{document}
